\documentclass[UTF8,12pt]{ctexart}
\usepackage[a4paper,margin=2.05cm]{geometry}
\usepackage{amsmath,amssymb}
\usepackage{array}
\usepackage{enumitem}
\usepackage{fancyhdr}
\usepackage{lastpage}
\usepackage[colorlinks,
linkcolor=black,
anchorcolor=black,
citecolor=black
]{hyperref}

\setlength{\parindent}{2em}
\setlength{\parskip}{0.2em}
\linespread{1.08}
\pagestyle{fancy}
\fancyhf{}
\cfoot{第 \thepage\ 页，共 \pageref{LastPage} 页}
\renewcommand{\headrulewidth}{0pt}

\newcommand{\blank}[1][4em]{\underline{\hspace{#1}}}
\newcommand{\examsection}[1]{\par\addvspace{0.7em}\noindent{\large\bfseries #1}\par\addvspace{0.25em}}
\newcommand{\choice}[4]{%
  \begin{tabular}{@{}p{0.48\linewidth}p{0.48\linewidth}@{}}
  (A) #1 & (B) #2\\[0.4em]
  (C) #3 & (D) #4
  \end{tabular}
}

\begin{document}

\begin{center}
  {\LARGE\bfseries 《高等数学 A（下）》期中考试试题}\\[0.5em]
  {\large 2025--2026 学年 第二学期 / 福建师范大学数学与统计学院}\\[0.3em]
  排版：\href{https://github.com/Xuuyuan}{@许愿} -- \href{https://wefjnu.nekoark.com}{WeFJNU}
\end{center}

\examsection{一、单选题（每题 3 分，共 15 分）}
\begin{enumerate}[label={\arabic*、},leftmargin=2em,itemsep=0.55em]
\item 下列函数组在其定义区间内线性无关的是（\quad）。

\choice{$\ln x,\ \ln x^2$；}
{$\sin x,\ \cos x$；}
{$\sin x\cos x,\ \sin 2x$；}
{$e^x,\ 2e^x$。}

\item 二次曲面 $z=\dfrac{x^2}{4}-\dfrac{y^2}{9}+z^2$ 与平面 $z=2$ 相截，其截痕是平面 $z=2$ 上的（\quad）。

\choice{双曲线；}{抛物线；}{椭圆；}{直线。}

\item 空间曲线
$\begin{cases}2x-y^2-z^2=0,\\ z=1\end{cases}$
在 $xOy$ 面上的投影曲线方程为（\quad）。

\begin{tabular}{@{}p{0.48\linewidth}p{0.48\linewidth}@{}}
(A) $\begin{cases}2x-y^2=1,\\x=0;\end{cases}$ &
(B) $\begin{cases}2x-y^2=1,\\y=0;\end{cases}$ \\[1.4em]
(C) $\begin{cases}2x-y^2=1,\\z=1;\end{cases}$ &
(D) $\begin{cases}2x-y^2=1,\\z=0.\end{cases}$
\end{tabular}

\item 设 $\varphi_1(x),\varphi_2(x)$ 为二阶非齐次线性微分方程
$y''+p(x)y'+q(x)y=f(x)$ 的二个不同的非零解，则下述结论正确的是（\quad）。

\choice{$\varphi_1(x)-\varphi_2(x)$ 是该方程的解；}
{$\varphi_1(x)+\varphi_2(x)$ 是该方程的解；}
{$\varphi_1(x)-\varphi_2(x)$ 是对应齐次方程的解；}
{$\varphi_1(x)+\varphi_2(x)$ 是对应齐次方程的解。}

\item 对于二元函数 $z=f(x,y)$，下述结论正确的是（\quad）。

\choice{全微分存在，则偏导数必连续；}
{全微分存在，则偏导数不一定存在；}
{偏导数连续，则全微分存在；}
{偏导数不连续，则全微分不存在。}
\end{enumerate}

\examsection{二、填空题（每小题 3 分，共 15 分）}
\begin{enumerate}[label={\arabic*、},leftmargin=2em,itemsep=0.55em]
\item $\displaystyle \lim_{(x,y)\to(0,0)}\frac{1-\cos(xy)}{x^2+y^2}=$ \blank[8em]。

\item 曲线 $r=2\cos\theta$，$\theta\in\left[-\dfrac{\pi}{2},\dfrac{\pi}{2}\right]$ 的弧长为 \blank[8em]。

\item 点 $(1,0,1)$ 到平面 $x-2y-z+1=0$ 的距离为 \blank[8em]。

\item 平面 $x+2y-2z+3=0$ 与平面 $x-y+z-3=0$ 的夹角余弦值为 \blank[8em]。

\item 设 $z=e^x\ln y$，则
$\displaystyle \left.\frac{\partial^5 z}{\partial x^2\partial y^3}\right|_{(0,1)}=$ \blank[8em]。
\end{enumerate}

\examsection{三、（8 分）}
求由曲线 $y=\sin x$ 和 $y=\cos x$ 以及直线 $x=0$ 和 $x=\dfrac{\pi}{2}$ 所围成图形面积。

\examsection{四、（8 分）}
求过点 $(2,0,-3)$ 且与直线
$\begin{cases}x-2y+4z=7,\\3x+5y-2z=-1\end{cases}$
垂直的平面方程。

\examsection{五、（8 分）}
设函数 $\varphi(x)$ 连续且满足方程
$\displaystyle \varphi(x)=\sin x+\int_0^x\varphi(t)\,\mathrm{d}t$，
求 $\varphi(x)$。

\examsection{六、（8 分）}
求微分方程
$\displaystyle \frac{\mathrm{d}^2y}{\mathrm{d}x^2}
-2\frac{\mathrm{d}y}{\mathrm{d}x}-3y=3x+2$
的通解。

\examsection{七、（8 分）}
设 $z=e^{xy}+f(x-y)$，其中 $f$ 为可微函数，求 $\mathrm{d}z$。

\examsection{八、（10 分）}
求解初值问题
\[
\begin{cases}
\dfrac{\mathrm{d}y}{\mathrm{d}x}=\dfrac{x}{y}+\dfrac{y}{x},\\[0.4em]
y(1)=2.
\end{cases}
\]

\examsection{九、（12 分）}
设函数
\[
f(x,y)=
\begin{cases}
\dfrac{x^3-y^3}{x^2+y^2}, & (x,y)\ne(0,0),\\[0.6em]
0, & (x,y)=(0,0),
\end{cases}
\]
判断该函数在 $(0,0)$ 处是否连续，偏导数是否存在，以及是否可微，并说明理由。

\examsection{十、（8 分）}
已知曲线 $\Gamma$ 为球面 $x^2+y^2+z^2=5$ 和平面 $z=1$ 的交线，求 $\Gamma$ 绕 $y$ 轴旋转一周所得到的曲面方程。

\end{document}
